\documentclass[10pt,a4paper]{amsart}
\usepackage{amsmath,amssymb,amsthm,graphicx}
% Retained minimal preamble for the Stacks Project chapter more-algebra.tex (More on Algebra).
% The chapter file pinned-original-source.tex carries only \input{preamble} on its first line, so no
% preamble text lives inside the chapter itself.  The Stacks Project uses one single shared preamble for
% every chapter; the snapshot stages it identically as sources/stacks-chow/preamble.tex and
% sources/stacks-etale-cohomology/preamble.tex.  The block below is copied mechanically, line for line,
% from sources/stacks-chow/preamble.tex lines 14-18 (the Xy-pic setup) and 154-217 (the native theorem
% declarations together with the mathematical macros).  Nothing was added, removed or redefined; the
% wrapper class and the amsmath/amssymb/amsthm/graphicx base packages are supplied by the assembler.
% Chapter: more-algebra.tex -- More on Algebra.
% For commutative diagrams we use Xy-pic
\usepackage[all]{xy}

% We use 2cell for 2-commutative diagrams.
\xyoption{2cell}

% Theorem environments.
%
\theoremstyle{plain}
\newtheorem{theorem}[subsection]{Theorem}
\newtheorem{proposition}[subsection]{Proposition}
\newtheorem{lemma}[subsection]{Lemma}

\theoremstyle{definition}
\newtheorem{definition}[subsection]{Definition}
\newtheorem{example}[subsection]{Example}
\newtheorem{exercise}[subsection]{Exercise}
\newtheorem{situation}[subsection]{Situation}

\theoremstyle{remark}
\newtheorem{remark}[subsection]{Remark}
\newtheorem{remarks}[subsection]{Remarks}

\numberwithin{equation}{subsection}

% Macros
%
\def\lim{\mathop{\mathrm{lim}}\nolimits}
\def\colim{\mathop{\mathrm{colim}}\nolimits}
\def\Spec{\mathop{\mathrm{Spec}}}
\def\Hom{\mathop{\mathrm{Hom}}\nolimits}
\def\Ext{\mathop{\mathrm{Ext}}\nolimits}
\def\SheafHom{\mathop{\mathcal{H}\!\mathit{om}}\nolimits}
\def\SheafExt{\mathop{\mathcal{E}\!\mathit{xt}}\nolimits}
\def\Sch{\mathit{Sch}}
\def\Mor{\mathop{\mathrm{Mor}}\nolimits}
\def\Ob{\mathop{\mathrm{Ob}}\nolimits}
\def\Sh{\mathop{\mathit{Sh}}\nolimits}
\def\NL{\mathop{N\!L}\nolimits}
\def\CH{\mathop{\mathrm{CH}}\nolimits}
\def\proetale{{pro\text{-}\acute{e}tale}}
\def\etale{{\acute{e}tale}}
\def\QCoh{\mathit{QCoh}}
\def\Ker{\mathop{\mathrm{Ker}}}
\def\Im{\mathop{\mathrm{Im}}}
\def\Coker{\mathop{\mathrm{Coker}}}
\def\Coim{\mathop{\mathrm{Coim}}}

% Boxtimes
%
\DeclareMathSymbol{\boxtimes}{\mathbin}{AMSa}{"02}

%
% Macros for moduli stacks/spaces
%
\def\QCohstack{\mathcal{QC}\!\mathit{oh}}
\def\Cohstack{\mathcal{C}\!\mathit{oh}}
\def\Spacesstack{\mathcal{S}\!\mathit{paces}}
\def\Quotfunctor{\mathrm{Quot}}
\def\Hilbfunctor{\mathrm{Hilb}}
\def\Curvesstack{\mathcal{C}\!\mathit{urves}}
\def\Polarizedstack{\mathcal{P}\!\mathit{olarized}}
\def\Complexesstack{\mathcal{C}\!\mathit{omplexes}}
% \Pic is the operator that assigns to X its picard group, usage \Pic(X)
% \Picardstack_{X/B} denotes the Picard stack of X over B
% \Picardfunctor_{X/B} denotes the Picard functor of X over B
\def\Pic{\mathop{\mathrm{Pic}}\nolimits}
\def\Picardstack{\mathcal{P}\!\mathit{ic}}
\def\Picardfunctor{\mathrm{Pic}}
\def\Deformationcategory{\mathcal{D}\!\mathit{ef}}

% --- appended 2026-10-12 (Chief): the project own definitions of the reference/slogan/history
% display environments, copied mechanically from the shared Stacks preamble lines 8-12
% (\usepackage{verbatim} plus the three \newenvironment lines, which the chapter preamble keeps
% verbatim from the project preamble). Nothing added, removed or redefined: these environments
% carry no mathematical content and the project itself typesets them as comments. Without these
% three lines a counted statement whose own text uses them fails with Environment slogan undefined.
\usepackage{verbatim}
\newenvironment{reference}{\comment}{\endcomment}
\newenvironment{slogan}{\comment}{\endcomment}
\newenvironment{history}{\comment}{\endcomment}
\usepackage{cleveref}
\crefname{originallemma}{original lemma}{original lemmas}
\Crefname{originallemma}{Original lemma}{Original lemmas}
\makeatletter
\newlabel{algebra-definition-quasi-finite}{{00PL}{source}{}{sourcepointer-1}{}}
\newlabel{algebra-definition-quasi-finite@cref}{{[originallemma][0][]00PL}{[1][0][]source}{}{}{}}
\makeatother
\crefname{originallemma}{original lemma}{original lemmas}
\Crefname{originallemma}{Original lemma}{Original lemmas}
\makeatletter
\newlabel{algebra-lemma-etale-quasi-finite}{{00U5}{source}{}{sourcepointer-2}{}}
\newlabel{algebra-lemma-etale-quasi-finite@cref}{{[originallemma][0][]00U5}{[1][0][]source}{}{}{}}
\makeatother
\crefname{originallemma}{original lemma}{original lemmas}
\Crefname{originallemma}{Original lemma}{Original lemmas}
\makeatletter
\newlabel{algebra-lemma-integral-no-inclusion}{{00GT}{source}{}{sourcepointer-3}{}}
\newlabel{algebra-lemma-integral-no-inclusion@cref}{{[originallemma][0][]00GT}{[1][0][]source}{}{}{}}
\makeatother
\crefname{originallemma}{original lemma}{original lemmas}
\Crefname{originallemma}{Original lemma}{Original lemmas}
\makeatletter
\newlabel{algebra-lemma-local-flat-ff}{{00HR}{source}{}{sourcepointer-4}{}}
\newlabel{algebra-lemma-local-flat-ff@cref}{{[originallemma][0][]00HR}{[1][0][]source}{}{}{}}
\makeatother
\crefname{originallemma}{original lemma}{original lemmas}
\Crefname{originallemma}{Original lemma}{Original lemmas}
\makeatletter
\newlabel{algebra-lemma-minimal-prime-image-minimal-prime}{{0CAN}{source}{}{sourcepointer-5}{}}
\newlabel{algebra-lemma-minimal-prime-image-minimal-prime@cref}{{[originallemma][0][]0CAN}{[1][0][]source}{}{}{}}
\makeatother
\crefname{originallemma}{original lemma}{original lemmas}
\Crefname{originallemma}{Original lemma}{Original lemmas}
\makeatletter
\newlabel{algebra-lemma-surjective-on-spec-units}{{0B7C}{source}{}{sourcepointer-6}{}}
\newlabel{algebra-lemma-surjective-on-spec-units@cref}{{[originallemma][0][]0B7C}{[1][0][]source}{}{}{}}
\makeatother
\crefname{originallemma}{original lemma}{original lemmas}
\Crefname{originallemma}{Original lemma}{Original lemmas}
\makeatletter
\newlabel{algebra-lemma-unramified-quasi-finite}{{02UR}{source}{}{sourcepointer-7}{}}
\newlabel{algebra-lemma-unramified-quasi-finite@cref}{{[originallemma][0][]02UR}{[1][0][]source}{}{}{}}
\makeatother
\title{For a flat localized etale, unramified, quasi-finite, integral or generic-fibre-unique map of domains the contraction of any nonzero ideal is nonzero}
\author{The Stacks Project Authors}
\date{stacks-project:more-algebra.tex:local-snapshot-20261012; GNU FDL 1.2 or later}
\begin{document}
\maketitle
\setcounter{subsection}{1}
\noindent\textit{Editorial note.} All delivered mathematical text is the author's own, copied from the pinned Stacks Project chapter file more-algebra.tex; no statement or proof step was rewritten, reordered or removed.\par\medskip
\noindent\textit{Editorial note.} The chapter file pinned-original-source.tex carries only \textbackslash{}input\{preamble\} on its first line, so no preamble text lives inside the chapter. The retained file sources/stacks-more-algebra/retained-preamble.tex is a mechanical, line-for-line copy of the project's single shared preamble as staged in the snapshot; nothing was added, removed or redefined and no package was added.\par\medskip
\noindent\textit{Editorial note.} The source URL is the source-level Stacks Project URL, as required when more than one card is produced from one source; the per-claim tag is carried by claim\_key.\par\medskip
\noindent\textit{Editorial note.} Seven enumerated hypotheses appear in the counted statement span and the proof treats them in the order case (7) first and then the reduction from the others; that order is the author's and is reproduced. The seven references inside the counted proof span are recorded as non-load-bearing navigation pointers; they include cross-chapter tags and same-chapter label lines, and the applications that consume them stay inside the counted proof span.\par\medskip
\noindent\textit{Allowed standard prerequisites (editorial).}\begin{itemize}
\item Fibres of Spec(B) -> Spec(A) and the generic point, and the behaviour of the generic point under a flat injective map. used at proof lines 31723-31729 and 31743-31748. Is the geometry that lets case (7) be identified as the weakest hypothesis.
\item Flat local homomorphisms are faithfully flat, so a flat map of domains is injective. used at proof lines 31740-31742. Guarantees that the generic point of Spec(B) maps to the generic point of Spec(A).
\item The tensor product with the fraction field as a filtered colimit of localizations. used at proof lines 31733-31734. Converts the unit statement over B tensor\_A K into a unit statement over some B\_f, which is what produces the nonzero element of A cap J.
\item No nontrivial specializations between fibre points for integral, quasi-finite, unramified and etale maps. used at proof lines 31749-31759. Reduces the remaining hypotheses to case (7).
\end{itemize}\medskip
\ifdefined\hypertarget\hypertarget{sourcepointer-1}{}\fi
\noindent\textit{Source reference.} Original lemma 00PL: focused-seeds/AGP/sources/stacks-snapshot-20261011/tags.txt, source lines 935--935. Cross-chapter navigation pointer into the Stacks Project; the standard fact it supplies is applied inside the counted proof span and the argument resting on it stays there.\par\medskip
\ifdefined\hypertarget\hypertarget{sourcepointer-2}{}\fi
\noindent\textit{Source reference.} Original lemma 00U5: focused-seeds/AGP/sources/stacks-snapshot-20261011/tags.txt, source lines 1099--1099. Cross-chapter navigation pointer into the Stacks Project; the standard fact it supplies is applied inside the counted proof span and the argument resting on it stays there.\par\medskip
\ifdefined\hypertarget\hypertarget{sourcepointer-3}{}\fi
\noindent\textit{Source reference.} Original lemma 00GT: focused-seeds/AGP/sources/stacks-snapshot-20261011/tags.txt, source lines 619--619. Cross-chapter navigation pointer into the Stacks Project; the standard fact it supplies is applied inside the counted proof span and the argument resting on it stays there.\par\medskip
\ifdefined\hypertarget\hypertarget{sourcepointer-4}{}\fi
\noindent\textit{Source reference.} Original lemma 00HR: focused-seeds/AGP/sources/stacks-snapshot-20261011/tags.txt, source lines 653--653. Cross-chapter navigation pointer into the Stacks Project; the standard fact it supplies is applied inside the counted proof span and the argument resting on it stays there.\par\medskip
\ifdefined\hypertarget\hypertarget{sourcepointer-5}{}\fi
\noindent\textit{Source reference.} Original lemma 0CAN: focused-seeds/AGP/sources/stacks-snapshot-20261011/tags.txt, source lines 15443--15443. Cross-chapter navigation pointer into the Stacks Project; the standard fact it supplies is applied inside the counted proof span and the argument resting on it stays there.\par\medskip
\ifdefined\hypertarget\hypertarget{sourcepointer-6}{}\fi
\noindent\textit{Source reference.} Original lemma 0B7C: focused-seeds/AGP/sources/stacks-snapshot-20261011/tags.txt, source lines 14100--14100. Cross-chapter navigation pointer into the Stacks Project; the standard fact it supplies is applied inside the counted proof span and the argument resting on it stays there.\par\medskip
\ifdefined\hypertarget\hypertarget{sourcepointer-7}{}\fi
\noindent\textit{Source reference.} Original lemma 02UR: focused-seeds/AGP/sources/stacks-snapshot-20261011/tags.txt, source lines 3742--3742. Cross-chapter navigation pointer into the Stacks Project; the standard fact it supplies is applied inside the counted proof span and the argument resting on it stays there.\par\medskip
\begin{lemma}
\label{lemma-ideal-inverse-nonzero}
Let $A$ and $B$ be domains and let $A \to B$ be a ring map.
Assume $A \to B$ has additionally at least one of the following properties
\begin{enumerate}
\item it is the localization of an \'etale ring map,
\item it is flat and the localization of an unramified ring map,
\item it is flat and the localization of a quasi-finite ring map,
\item it is flat and the localization of an integral ring map,
\item it is flat and there are no nontrivial specializations between points
of fibres of $\Spec(B) \to \Spec(A)$,
\item $\Spec(B) \to \Spec(A)$ maps the generic point to the generic
point and there are no nontrivial specializations between points of fibres, or
\item exactly one point of $\Spec(B)$ is mapped to the generic
point of $\Spec(A)$.
\end{enumerate}
Then $A \cap J$ is nonzero for every nonzero ideal $J$ of $B$.
\end{lemma}

\begin{proof}
Proof in case (7).
Let $K$, resp.\ $L$ be the fraction field of $A$, resp.\ $B$.
By Algebra, Lemma
\ref{algebra-lemma-minimal-prime-image-minimal-prime} we see that
the unique point of $\Spec(B)$ which maps to the generic
point $(0) \in \Spec(A)$ is $(0) \in \Spec(B)$.
We conclude that $B \otimes_A K$ is a ring with a unique prime
ideal whose residue field is $L$ (in fact it is equal to $L$ but
we do not need this).
Choose $b \in J$ nonzero. Then $b$ maps to a unit of $L$.
Hence $b$ maps to a unit of $B \otimes_A K$
(Algebra, Lemma \ref{algebra-lemma-surjective-on-spec-units}).
Since $B \otimes_A K = \colim_{f \in A \setminus \{0\}} B_f$
we see that $b$ maps to a unit of $B_f$ for some $f \in A$ nonzero.
This means that $b b' = f^n$ for some $b' \in B$ and $n \geq 1$.
Thus $f^n \in A \cap J$ as desired.

\medskip\noindent
In the rest of the proof, we show that each of the other assumptions imply (7).
Under assumptions (1) -- (5), the ring map $A \to B$ is flat and
hence $A \to B$ is injective (since flat local homomorphisms are
faithfully flat by Algebra, Lemma \ref{algebra-lemma-local-flat-ff}).
Hence the generic point of $\Spec(B)$
maps to the generic point of $\Spec(A)$. Now, if there are no
nontrivial specializations
between points of fibres of $\Spec(B) \to \Spec(A)$, then of course
this generic point of $\Spec(B)$ has to be the unique point mapping
to the generic point of $\Spec(A)$. So (6) implies (7).
Finally, to finish we show that in cases (1) -- (5) there are no
nontrivial specializations
between the points of fibres of $\Spec(B) \to \Spec(A)$.
Namely, see
Algebra, Lemma \ref{algebra-lemma-integral-no-inclusion} for
the integral case,
Algebra, Definition \ref{algebra-definition-quasi-finite}
for the quasi-finite case, and use that
unramified and \'etale ring maps are quasi-finite
(Algebra, Lemmas \ref{algebra-lemma-unramified-quasi-finite} and
\ref{algebra-lemma-etale-quasi-finite}).
\end{proof}

\end{document}
